\section{Quotient Rings }
\begin{theorem}
If $I$ is an ideal of a ring $R$ the cosets $x + I$ form a ring under the sum $ (x + I) + (y + I) = (x+y) + I$ and products
$(x + I ) (y + I) = xy + I$.
\end{theorem}

The ring of cosets is called the {\twelveit Quotient Ring of I in R} and it is written $R/I$.  The mapping $a \goes a' = a + I$ for $a \in R$, which carries each element of $R$ into the coset 
containing it is a homomorphism of $R$ into the quotient ring $R/I$ and the elements mapped on zero in this homomorphism are exactly the elements of the ideal $I$. In the homomorphic ring, the zero is the coset $0 + I$, hence the elements of $I$ are mapped onto the zero coset.  

\begin{theorem}
If a homomorphism $H$ maps $R$ on $R'$ and exactly the elements of $I$ on zero, then $R'$ is isomorphic to the quotient ring $R/I$. 
\end{theorem}